\subsection{Simultaneously primitive bundle annuli}
\label{marked-cut:primitive-annuli}

The next extension handles nondisk bases when at least one boundary component
is not fully attached to the core. It uses the topology of the embedded
interval bundle to certify a simultaneous primitive family. It does not infer
primitivity from an abelian word, or merely from the topology of the annuli.
The certificate concludes the genus census of the whole cut; it supplies no
explicit transported boundary pattern for a later hierarchy stage.

A primitive family of annuli in a handlebody has pairwise disjoint dual disks,
each meeting its designated annulus in one spanning arc and avoiding all the
others. This is the standard geometric notion used in
Schultens--Weidmann~\cite[Definition~6 and Lemma~2]{schultens-weidmann-primitive}.
Their stated gluing lemma assumes essential annuli on the opposite side.
The special source-bundle argument below gives an explicit handle cancellation
proof and does not require a test of that additional assumption.

\begin{proposition}[A spare base boundary gives a primitive family]
\label{marked-cut:primitive-family}
Let $P\to F$ be an orientable interval-bundle component over a connected
compact surface with $b>0$ boundary components. Suppose $a$ distinct complete
base boundary circles are core-labelled, with $0<a<b$. Their vertical annuli
form a simultaneous primitive family in the handlebody $P$, relative to all
other proper interval attachments. Cutting along its dual disks leaves a
connected handlebody of genus $1-\chi(F)-a$. The assertion allows
nonorientable $F$ and a twisted interval bundle.
\end{proposition}
\begin{proof}
Choose a boundary component $C_0$ outside the $a$ complete attaching circles.
It has an open arc outside the core-labelled set: otherwise the complete
ordered boundary skeleton would have counted it among those circles. An
unported boundary component is entirely outside that set.

In $F$, choose disjoint properly embedded arcs $\gamma_i$ joining each
fully attached circle $C_i$ to distinct points in this free arc of $C_0$.
One way to see their existence is to join distinct boundary components
successively. An arc between two distinct boundary components is
nonseparating; cutting along it merges those components and leaves the
surface connected. The subsequent arcs can meet the remaining free portion
of $C_0$. Reversing the cuts gives the required disjoint arcs. This argument
uses neither orientability nor a bound on the genus.

Each preimage $D_i=p^{-1}(\gamma_i)$ is a properly embedded disk, since an
interval bundle over an arc is a rectangle. Its boundary meets the designated
vertical annulus in a single fibre, a spanning arc, and is disjoint from
all other complete attaching annuli. Its other vertical end lies in the
chosen free part of $C_0$; its horizontal edges lie on the horizontal bundle
boundary. Thus the disks avoid every remaining proper interval attachment.
Their interiors are disjoint because the base arcs are disjoint.

Cutting $P$ along these disks is the interval bundle over the surface cut
along the $\gamma_i$. This base is connected, bordered and has Euler
characteristic $\chi(F)+a$. Its orientable total interval bundle is a
handlebody of genus $1-\chi(F)-a$, by the bordered-bundle argument already
used above. Reattaching the disks supplies $a$ distinguished one-handles,
each traversed exactly once by its corresponding annulus and avoided by
the other attaching annuli and disks.
\end{proof}

\begin{proposition}[Whole-cut gluing with primitive annuli]
\label{marked-cut:primitive-census}
First absorb disk-base full circles and certify the resulting core
components as handlebodies. Suppose every other ported bundle has a bordered
base and its full-circle attachments satisfy $a<b$. Then the remaining
annular and disk interfaces glue these pieces into handlebodies. In each
connected attachment graph component the genus is
\[
 g=\sum_{v\in V}(g_v-a_v)+|E|-|V|+1,
\]
where $g_v$ is the core or bundle genus, $a_v$ counts primitive annuli at a
bundle vertex and is zero at core vertices, and $E$ includes both disk and
primitive-annulus attachments. Unported bundle components retain their
separate binary multiplicities.
\end{proposition}
\begin{proof}
For an annulus gluing, cut the annulus along a spanning arc to a disk.
Gluing that disk is represented by a connecting one-handle; restoring the
annulus is represented by a two-handle whose attaching curve runs once
around each of the two annuli, with connecting arcs through that handle.
This is the usual relative handle description obtained by thickening a
one-vertex, one-edge cell structure on the annulus. The connecting handle
joins the two corresponding components, or creates a cycle if they are
already connected.

Use the distinguished one-handle on the bundle side supplied by
Proposition~\ref{marked-cut:primitive-family}. The new two-handle intersects
its cocore exactly once. No other attachment intersects that cocore, because
the dual disks avoid all other annuli and proper interval interfaces.
The pair therefore cancels geometrically. All these cancellations may be
performed successively: each annulus uses its own distinguished handle.
The opposite core annulus may be compressible; it adds no intersection with
a cocore lying in this bundle. An arbitrary gluing map of an annulus still
carries its essential core circle once around the target annulus, so a twist
of the gluing does not change the unique bundle-side intersection.

What remains is a family of handlebodies of genera $g_v-a_v$, joined along
one disk for every original disk or primitive annulus interface. A spanning
tree of the attachment graph joins the pieces without adding genus; each
remaining edge adds one one-handle. Orientability is inherited from the
source cut. This proves the formula, including
repeated attachments between the same two pieces. Bundle-side dual disks
avoid the original proper disk interfaces, so this cancellation does not
invalidate their simultaneous use.

For an independent Euler check, write $D$ and $A$ for the disk and annular
edge counts. Since $|E|=D+A$ and $\sum a_v=A$, the formula gives
$\chi=\sum_v(1-g_v)-D$. This is precisely the source gluing identity:
disks have Euler characteristic one and annuli zero. Previously absorbed
disk bundles have already been included in the augmented core.
\end{proof}

\paragraph{A finite certificate rather than an expanded disk system.}
The producer groups the complete attaching circles by their authenticated
base owner. Disk bases are left to the source-chart absorption adapter.
Every other such base must have multiplicity one, $b>a$, and
$1-\chi(F)\ge a$. The record names all its annular attachment indices,
$b-a$ spare boundary components, and reduced genus $1-\chi(F)-a$.
The boundary checker has already proved which circles are wholly attached
and which maximal attachments are proper intervals. Counting a proper
interval as a complete annulus would change that source-verified record.

The independent module \path{bundle_primitive_annuli_verify} reconstructs
owners, complete-circle lists, spare counts and genus deductions without
calling the producer. Only these independently established base owners may
be skipped by the disk-absorption checker. The latter still authenticates
every actually absorbed two-handle. Core cancellation then proceeds as
before; the final graph checker uses the reduced bundle genera and includes
the primitive annuli as connecting edges. Producer and checker use separate
graph algorithms, and both retain the original source-cut Euler identity.

These proofs use schema \path{normal-cut-handlebodies-v3}. Schemas v1 and v2
remain replayable, and a cut without primitive annuli retains its earlier
schema. A schema downgrade cannot omit the new argument: the earlier
checkers still reject unabsorbed annuli. The six weighted queries, shared
orbit budgets and source input contract are unchanged. A refused family or
stalled augmented core proves neither non-handlebodyness nor knottedness.

\paragraph{Size, running time and remaining obligations.}
There are at most $P$ complete attaching circles. The new record has $O(P)$
indices and binary integers; no dual disk system, base genus, normal sheet
stack or boundary component family is expanded. Grouping uses $O(P)$ finite
incidences and sorting at most $O(P\log P)$ comparisons. If the source profile
integers have at most $B$ bits, the added arithmetic and record sizes are
polynomial in $T+B$, with record size $O(P(B+\log(T+P)))$ bits. The previous
source-orbit and saturated-core bounds therefore still dominate in order.

The spare-boundary hypothesis is sufficient, not necessary. A M\"obius band
attached along its only boundary and an annulus attached along both boundaries
remain outside this argument. Explicit transport of arbitrary boundary
patterns, handling every inherited core, useful-surface coverage and a global
quasi-polynomial hierarchy bound remain separate requirements.

\paragraph{Source-cut audit and prior successes.}
The refreshed 106-case audit certifies 78 cuts and independently replays all
78 proofs. Every one of the previous 63 successes is retained, and the 15
additional successes are exactly the eligible cases in the preceding profile
diagnostic. Regina constructs the actual cut and recognizes its connected
components in 99 cases, with exact agreement on all 71 native successes below
the literal expansion cap. Seven large-coordinate successes use independent
source replay without expansion; there are no oracle timeouts. Records and
runtime pins are retained in \path{data/primitive-annuli-audit.json}.

The 28 remaining refusals are 11 unsupported annular interfaces, seven failed
augmented core cancellations (including three square leaf obstructions), and
ten cores with three-handles. They carry no negative recognition conclusion.
The new nonorientable-base example certifies one ball and two genus-three
handlebodies; the five-tetrahedron multiple-annulus example certifies one ball
and one genus-two handlebody. Both match independently constructed Regina cuts.

The 22 focused tests pass. They cover simultaneous annuli, repeated attachment
graph edges, a nonorientable base, source relabellings, no-spare-boundary
refusals, forged counts and indices, schema downgrades, disabled producers,
hexadecimal transport, shared budget caps, multiplicity $2^{1024}$, and legacy
v1/v2 replay. All low-level family arguments remain conditional on the source
proofs; the whole-cut verifier checks those sources first.

